\documentclass[12pt,a4paper]{article}
\usepackage[T2A]{fontenc}
\usepackage[utf8]{inputenc}
\usepackage[russian]{babel}
\usepackage{geometry}
\usepackage{amsmath,amssymb,mathtools}
\usepackage{array,booktabs}
\usepackage{xcolor}
\usepackage{tikz}
\usepackage{hyperref}
\geometry{left=23mm,right=23mm,top=22mm,bottom=24mm}
\hypersetup{colorlinks=true,linkcolor=blue!55!black,urlcolor=blue!55!black}
\setlength{\parindent}{0pt}
\setlength{\parskip}{4.5pt}
\renewcommand{\arraystretch}{1.25}
\newcommand{\rubric}[1]{\par\smallskip\noindent\textbf{#1}\par\nobreak\smallskip}
\begin{document}
\begin{titlepage}
\centering
\vspace*{0.28\textheight}
{\Huge\bfseries Ревью домашней работы\par}
\vspace{2.4cm}
{\large Автор: Лёвин Артём Александрович\par}
\vspace{0.8cm}
{\large 17 октября 2026 года\par}
\vfill
\begin{tikzpicture}
\draw[blue!35!black,line width=0.65pt] (0,0) .. controls (2.0,0.38) and (4.4,-0.38) .. (6.6,0);
\end{tikzpicture}
\vspace{0.8cm}
\end{titlepage}

\newpage
\section*{Сводная таблица замечаний}
\small
\begingroup\setlength{\tabcolsep}{4pt}
\begin{tabular}{@{}p{1.1cm}p{3.2cm}p{4.4cm}p{3.0cm}p{2.7cm}@{}}
\toprule
\textbf{№} & \textbf{Тема} & \textbf{Суть недоработки} & \textbf{Ваш ответ} & \textbf{Правильный ответ}\\
\midrule
\hyperref[task:7]{7} & Сложные проценты; поиск ставки & Найден коэффициент роста, но не вычислена ставка в процентах. & $k=1{,}1$; ставка не указана. & $10\%$ годовых\\
\bottomrule
\end{tabular}
\endgroup
\normalsize

\newpage
\vspace*{\fill}
\begin{center}
\section*{Разбор ошибок}
\end{center}
\vspace*{\fill}

\newpage
\phantomsection\label{task:7}
\section*{Задание № 7}

\rubric{Текст задания}
Через 3 года вклад 500 тыс. руб. без дополнительных операций вырос до 665,5 тыс. руб. Найдите постоянную годовую процентную ставку.

\rubric{Ваше решение}
Обозначен коэффициент ежегодного увеличения вклада через $k$. Записаны и решены уравнения
\[
665{,}5=500\cdot k^3,\qquad
k^3=1{,}331,\qquad
k=1{,}1.
\]

\rubric{Ваш ответ}
Записано $k=1{,}1$. Значение годовой процентной ставки отдельно не указано.

\rubric{Ошибка}
\textcolor{red}{Ошибка:} решение не доведено до требуемой величины. Вы правильно нашли \emph{коэффициент роста} $k$, однако по условию необходимо найти \emph{процентную ставку} $p$.

\rubric{Где именно возникает недоработка}
Последний правильный шаг --- $k=1{,}1$. Сразу после него требуется перевести коэффициент роста в процентную ставку. Этот завершающий переход отсутствует; неверного вычисления в записанных шагах нет.

\rubric{Почему этого недостаточно}
Коэффициент роста и процентная ставка --- не одно и то же. Коэффициент $1{,}1$ показывает, во сколько раз увеличивается сумма за год: новая сумма составляет $1{,}1$ от суммы в начале года. Но вклад увеличивается \emph{не на} $1{,}1\%$, а на долю $1{,}1-1=0{,}1$ от исходной суммы за год.

Общий смысл связи между величинами такой: при годовой ставке $p\%$ вклад после начисления процентов умножается на
\[
k=1+\frac{p}{100}.
\]
Единица в этой формуле сохраняет всю сумму, которая уже была на счёте, а дробь $p/100$ добавляет долю начисленных процентов. Поэтому из найденного коэффициента нужно сначала вычесть единицу, а затем перевести полученную долю в проценты, умножив на сто.

\rubric{Ключевая идея}
\textcolor{blue!60!black}{Ключевая идея:} сначала по отношению конечной и начальной сумм найти годовой коэффициент роста, затем определить ставку по формуле
\[
p=100(k-1).
\]

\rubric{Исправление от последнего правильного шага}
Ваш результат $k=1{,}1$ сохраняем. Вычитаем единицу, чтобы оставить только годовой прирост, и выражаем его в процентах:
\[
\frac{p}{100}=1{,}1-1=0{,}1,\qquad
p=0{,}1\cdot100=10.
\]
Следовательно, искомая ставка составляет $10\%$ годовых.

\rubric{Полное правильное решение}
Пусть $p\%$ --- постоянная процентная ставка в год. Поскольку дополнительные внесения и снятия отсутствуют, каждый год вся имеющаяся сумма увеличивается в одно и то же число раз:
\[
1+\frac{p}{100}.
\]
За три года этот множитель применяется трижды, поэтому
\[
500\left(1+\frac{p}{100}\right)^3=665{,}5.
\]
Делим обе части на $500$, чтобы получить уравнение только для коэффициента:
\[
\left(1+\frac{p}{100}\right)^3
=\frac{665{,}5}{500}=1{,}331.
\]
Замечаем, что $1{,}1^3=1{,}331$. Так как ежегодный коэффициент роста положителен, получаем
\[
1+\frac{p}{100}=1{,}1.
\]
Вычитаем единицу из обеих частей:
\[
\frac{p}{100}=0{,}1.
\]
Умножаем на сто:
\[
p=10.
\]
Значит, годовая процентная ставка равна $10\%$.

\rubric{Проверка}
При ставке $10\%$ годовой коэффициент действительно равен $1{,}1$. Последовательно вычислим суммы в тысячах рублей:
\[
\begin{aligned}
\text{через первый год:}\quad &500\cdot1{,}1=550,\\
\text{через второй год:}\quad &550\cdot1{,}1=605,\\
\text{через третий год:}\quad &605\cdot1{,}1=665{,}5.
\end{aligned}
\]
Конечная сумма совпала с условием, поэтому найденная ставка верна.

\rubric{Итог}
\textcolor{teal!60!black}{Итог:} вычисления до шага $k=1{,}1$ выполнены верно. Недоставало лишь перехода от коэффициента роста к ставке. \textbf{Ответ: $10\%$ годовых.}

\rubric{Задача на закрепление}
Вклад 400 тыс. руб. без пополнений и снятий за два года вырос до 576 тыс. руб. Найдите постоянную годовую процентную ставку при ежегодной капитализации.

\vfill
\begin{center}
\footnotesize Лёвин Артём Александрович эксклюзивно для Матвея Горбачева
\end{center}
\end{document}
